\section{Introduction}
\label{sec:introduction}

For a Noetherian regular local domain $A$ with fraction field $F$,
Gersten's conjecture asserts exactness of the canonical augmented residue
complex
\begin{equation}\label{eq:gersten}
 0\longrightarrow K_n(A)\longrightarrow K_n(F)
 \longrightarrow\bigoplus_{\operatorname{ht}\mathfrak p=1}
 K_{n-1}(\kappa(\mathfrak p))\longrightarrow\cdots .
\end{equation}
Here $n\geq0$, the groups are integral Quillen $K$-groups, and the
residues come from localization and the codimension filtration.
The first assertion is generic injectivity: a class that vanishes after
passage to the fraction field must already vanish on $A$.
We construct a counterexample to this assertion.

\begin{theorem}\label{thm:main}
Let $V$ be the ring of integers in the unramified extension of $\Q_5$ of
degree eight, and put
\[
 A=\bigl(V[x,y]/(5+x^4+y^4)\bigr)_{(5,x,y)}.
\]
Then $A$ is a Noetherian regular local domain of dimension two and
\[
 \ker\bigl(K_5(A)\longrightarrow K_5(\Frac A)\bigr)\ne0.
\]
Consequently the unrestricted integral Gersten conjecture has a negative
answer.
\end{theorem}

The maximal ideal of $A$ is $(x,y)$, and $5\in(x,y)^4$.
Its special fiber has dimension one and embedding dimension two at the
closed point. Thus the regular total ring is ramified and is not smooth
over $V$. This distinction is central both to the example and to its
relation with positive results on Gersten's conjecture.

\paragraph{History and related results.}
Gersten's questions about transfer and local higher $K$-theory
\cite{Gersten1973Problems}, together with his work on
exact sequences \cite{Gersten1973Exact}, place the conjecture in the
study of localization and residue maps. Quillen formulated it for
arbitrary regular local rings and proved it for
semilocalizations of finite-type algebras over a field at regular primes
\cite[\S7, Conjecture~5.10 and Theorem~5.11]{Quillen1973Higher}.
Panin established the full equicharacteristic case
\cite[Theorem~A]{Panin2003}.

In mixed characteristic, Gillet and Levine's relative theorem reduces
the smooth local case to the associated discrete valuation ring case
\cite{GilletLevine1987}. Geisser and Levine prove finite-coefficient
Gersten exactness for semilocal rings essentially smooth over a discrete
valuation ring, for every finite coefficient order, including powers of
the residue characteristic \cite[Theorem~8.2]{GeisserLevine2000}; their
text credits B.~Kahn for pointing out this consequence. Skalit proved
finite-coefficient Gersten exactness for arbitrary unramified regular
local rings, again for every finite coefficient order
\cite[Corollary~5.7]{Skalit2017}. In mixed characteristic $(0,p)$,
unramified means $p\notin\mathfrak m^2$; our ring is excluded since
$5\in\mathfrak m^4$. Version~1 of
Druzhinin's preprint states a result for essentially smooth local
Henselian schemes over a discrete valuation ring
\cite[Theorem~4.5]{Druzhinin2025}. The smoothness hypotheses exclude the
structure map $V\to A$ of Theorem~\ref{thm:main}.

Feld's finite-coefficient example is a close methodological predecessor
\cite[Theorem~1.1 and Corollary~5.1]{Feld2026}. It gives a degree-two
augmentation kernel modulo three at a ramified regular local ring whose
integral $K_2$ map is injective. Its supported pushforward, character
averaging, cyclotomic scalar extension, conjugate splitting and
divisor-class detection also organize our construction. In that example,
an integral lift of the coefficient kernel has nonzero but
three-divisible generic image, as follows from integral injectivity and
the coefficient exact sequence. An integral generic kernel therefore
requires an additional argument. Here a vanishing result for integral
$K_4$ permits a cyclotomic coefficient class to lift to integral $K_5$.
Two Bott factors contribute the square of the cyclotomic character,
whose signs are realized by an elliptic quotient of the Fermat quartic.

Our companion paper, \emph{An integral degree-three Gersten counterexample}
\cite[Theorem~1.1]{OpenAIDegreeThree2026}, constructs an integral
$K_3$ kernel at a different ramified ring: the local ring at the special
point $[0:0:1]$ of the cyclotomic quintic model
$(X-Z)^5-XY^4+Y^5+\zeta Z^5=0$ over the integers of $\Q_5(\zeta)$,
where $\zeta$ is a primitive fifth root of unity.
Its proof uses a Steinberg relation and a norm valuation test on coherent
$K$-theory \cite[Lemma~4.2, Proposition~3.1, and \S5]{OpenAIDegreeThree2026}.
Here the global arithmetic lift, smooth Picard specialization criterion
and quartic construction are proved directly. Neither proof uses results
of the other paper, and neither determines a smallest degree of failure.

Our proof detects the reduction modulo five of an integral class already
killed over the fraction field. It therefore also gives generic
noninjectivity with coefficients $\Z/5$ for the same ring. This detection
does not determine the order of the integral class or whether its
rationalization is nonzero. The conclusion concerns the unrestricted
integral conjecture and does not settle its rational-coefficient variant.

\paragraph{Construction and proof guide.}
A class pushed forward from a closed divisor of $\Spec A$ vanishes over
$\Frac A$. We must construct such a class integrally and prove that its
pushforward is nonzero. These tasks account for the two branches of the
proof.

For the integral construction, put $S=\Z[\zeta,1/5]$.
Section~\ref{sec:coefficients} supplies coefficient multiplication and
its exact field isomorphisms at the low degrees used here.
Section~\ref{sec:arithmetic} chooses $\beta\in K_2(S;5)$, writes $b$
for left multiplication by $\beta$, and constructs
\[
 c\in K_5(S),\qquad c\bmod5=b^2[\zeta-1]_5,
\]
where $b^2(x)=\beta\cdot(\beta\cdot x)$ and $[v]_5$ is the coefficient
$K_1$-class of a unit. The obstruction
is $K_4(S)[5]$, the subgroup killed by five. Brauer reciprocity, geometric
coefficient detection and finite generation prove that this subgroup
vanishes. Conjugating the reduction of $c$ multiplies its two Bott
factors by the square of the cyclotomic character.

Write $H_0=\Frac V$, $L=H_0(\zeta)$, and let $\OO$ be the valuation
ring of $L$. Integral localization supplies the following route:
\[
 c\in K_5(S)\longmapsto c_L\in K_5(L)
 \longleftarrow\widetilde c\in K_5(\OO).
\]
The leftward arrow is restriction from $\OO$, and $\widetilde c$ is a
chosen preimage; its existence uses integral finite-field $K_4$ vanishing.
Section~\ref{sec:local-geometry} constructs a surjection $A\to\OO$.
Its closed immersion $i$ then gives the integral class
$\gamma=i_*\widetilde c\in K_5(A)$, with zero generic restriction.

For detection, flat extension from $V$ to $L$ splits this divisor into
four rational points on the generic Fermat quartic $X^4+Y^4=Z^4$.
Section~\ref{sec:specialization} proves a general criterion on an arbitrary
localization of a smooth affine curve. A hypothetical relation among
supported coefficient classes untwists, by two Bott operations, to
ordinary integral tame symbols. On a smooth proper arithmetic surface,
the coniveau identity $d_1^2=0$ makes their horizontal specialization a
principal divisor. A homomorphism from the entire special-fiber Picard
group can therefore test the proposed relation.

Section~\ref{sec:finite-geometry} constructs this test. The Fermat quartic
maps with degree two to the genus-one curve $w^2=s^4+1$. With a branch
point as origin, the four chosen specializations map to $Q$ or $-Q$,
with the same signs as the character square. Their weighted sum is
$4Q\ne0$ modulo five. Section~\ref{sec:main-proof} verifies all hypotheses
of the Picard criterion and proves that the coefficient reduction of
$\gamma$ remains nonzero after flat pullback. This detects the already
integral class $\gamma$.

Appendix~\ref{app:claimed-vanishing} treats a separate issue in
Mochizuki's version~8 preprint \cite{Mochizuki2020}. Its Theorem~5 would
annihilate the principal-divisor pushforward used here. We prove an
explicit failure of the strict functor assertion in Lemma~13(III), for
every nonzero regular parameter. That calculation concerns the supplied
proof construction, and by itself does not refute its theorem statement.
No argument in the kernel construction depends on the appendix.
